\chapter{Conclusion and Future Scope}

\section{Conclusion}
The HomeSOC project successfully culminated in the comprehensive design, programmatic development, verified rigorous testing, and eventual global cloud deployment of a fully functional, highly optimized lightweight Security Operations Center. Driven by the core objective to democratize cybersecurity visibility for localized residential networks, small-to-medium enterprise (SME) environments, and individual home-lab deployments, the fundamental engineering paradigm deliberately eschewed the exorbitant computational overhead and systemic complexity traditionally inherent within enterprise-grade Security Information and Event Management (SIEM) systems like Splunk or Elasticsearch.

Through strict adherence to an agile development testing module sequence, the architecture was explicitly segmented. The development of a Python-based endpoint host extraction agent successfully demonstrated that systemic metrics (CPU load, memory thresholds) and critical security logs (SSH access attempts, process invocations) could be autonomously scraped, mapped to serialized JSON formats, and transmitted reliably with negligible impact on the underlying execution node hardware resources. 

The successful implementation of the centralized routing Flask ingestion server, physically fortified specifically via SQLite Write-Ahead Logging (WAL) pragmas, categorically verified the structural throughput integrity of the proposed detection heuristics logic engine. The system unequivocally demonstrated that mathematical arithmetic string mapping alone—when applied accurately via modular database parameters—is singularly sufficient for flagging critical attack patterns natively without relying on heavy external third-party threat-feed API indexing services. 

Finally, the translation of encrypted database parameters into a highly aesthetic, fluid, zero-latency dark-mode web dashboard interface formally achieved the primary human-interaction goal. By actively synthesizing abstract analytical security telemetry into prominently visible, categorized, and interactive chronological "Toasts" detailing clear remedial actions, HomeSOC empowers the layperson administrator to actively engage, react, and neutralize network threats practically, fulfilling all outlined project objectives thoroughly.

\section{Future Scope and Advancements}
While the final deployment configuration (`deployment\_ready\_v7.zip`) perfectly encompasses the defined Minimum Viable Product (MVP) operational bounds, the modularity explicitly engineered into the core codebase enables significant scalable evolutionary pathways for future academic or enterprise research iterations.

\subsection{Machine Learning and Anomaly Intelligence}
Presently, the `detection.py` analytical logic is deterministically constrained by statically programmed mathematical threshold constants defined globally inside the `rules` SQL parameter structural block. While highly functional, this logic cannot adapt dynamically to slowly shifting baseline network profiles natively. Future architectures should explicitly substitute the simple Boolean execution matrices with supervised Machine Learning (ML) execution models—specifically utilizing lightweight Isolation Forest algorithms or pre-trained Recurrent Neural Networks (RNN) mapped via TensorFlow Lite string processing. This would allow the SOC engine specifically to actively "learn" normative data patterns continuously, automatically flagging abstract behavioral anomalies (e.g., a "living-off-the-land" powershell attack sequence) entirely devoid of explicitly mapped string syntax rules.

\subsection{API Authentication via Bearer Tokens}
To facilitate rapid local testing environments, the core ingestion routing `/api/logs` endpoint listener was engineered universally open, accepting unencrypted unauthenticated HTTP physical POST request injections natively from any loopback script routing interface. To securely scale the HomeSOC architecture across zero-trust configurations or exposed public IP domain topologies natively, explicit authentication protocols must inherently be fundamentally developed. Executing an OAuth2 or JSON Web Token (JWT) verification pipeline script explicitly requiring a cryptographic structural static Bearer Token Header dynamically authenticated per-node would categorically harden the application layer routing boundaries natively preventing external malicious remote log spoofing cyber attacks.

\subsection{Docker Containerization Topology}
Currently, the operational server infrastructure string is hosted directly natively executing via PythonAnywhere WSGI virtual environments structurally dynamically interpreting scripts explicitly via root file mapping directories configurations directly attached explicitly directly attached to the system interpreter arrays configurations parameters natively executing operations natively. Transforming the `app.py`, `detection.py`, and `db.sqlite` functional dependency ecosystems explicitly dynamically comprehensively completely into a strictly defined, self-contained `Dockerfile` image representation sequence array logic block explicitly completely containerizes the ecosystem natively entirely definitively. Utilizing physical Docker container volumes mapping explicitly allowing abstract scaling execution instances deployed safely across redundant Kubernetes orchestration clusters topologies natively inherently inherently definitively establishing total system high availability architectures faultlessly completely organically dynamically universally continuously safely inherently globally.

\begin{table}[h]
\centering
\small
\begin{tabularx}{\textwidth}{|X|X|X|}
\hline
\textbf{Proposed Enhancement Category} & \textbf{Current MVP Implementation} & \textbf{Future Architectural Paradigm Shift} \\ \hline
\textbf{Threat Heuristic Model} & Static Boolean Database SQL String Matching & Dynamic Supervised Machine Learning (RNN / Isolation Forest) \\ \hline
\textbf{Node Authentication Strategy} & Unverified open local network REST POSTs & Cryptographic JWT Bearer Token validation per request \\ \hline
\textbf{Infrastructure Topology} & Raw OS Python virtual environments & Containerised Docker logic spanning Kubernetes multi-cluster nodes \\ \hline
\textbf{Hardware Interaction Scope} & CPU, RAM, Disk, Auth log extraction & Unrestricted eBPF (Extended Berkeley Packet Filter) kernel-level interception \\ \hline
\end{tabularx}
\caption{Strategic summary outlining the planned evolutionary trajectory expanding the application capabilities.}
\label{tab:future_trajectory}
\end{table}

\section{Expanded Theoretical Framework: Future Implementations in Cybersecurity}
Machine learning (ML) is a field of study in artificial intelligence concerned with the development and study of statistical algorithms that can learn from data and generalize to unseen data, and thus perform tasks without explicit programming language instructions. Within a subdiscipline of machine learning, advances in the field of deep learning have allowed neural networks, a class of statistical algorithms, to surpass many previous machine learning approaches in performance.
Statistics and mathematical optimisation (mathematical programming) methods compose the foundations of machine learning. Data mining is a related field of study, focusing on exploratory data analysis (EDA) through unsupervised learning.
From a theoretical viewpoint, probably approximately correct learning provides a mathematical and statistical framework for describing machine learning. Most traditional machine learning and deep learning algorithms can be described as empirical risk minimisation under this framework.


== History ==

The term machine learning was coined in 1959 by Arthur Samuel, an IBM employee and pioneer in the field of computer gaming and artificial intelligence. The synonym self-teaching computers was also used during this time period.
The earliest machine learning program was introduced in the 1950s when Arthur Samuel invented a computer program that calculated the winning chance in checkers for each side, but the history of machine learning roots back to decades of human desire and effort to study human cognitive processes. In 1949, Canadian psychologist Donald Hebb published the book The Organization of Behavior, in which he introduced a theoretical neural structure formed by certain interactions among nerve cells. Hebb's model of neurons interacting with one another set a groundwork for how AIs and machine learning algorithms work under nodes, or artificial neurons used by computers to communicate data. Other researchers who have studied human cognitive systems contributed to the modern machine learning technologies as well, including logician Walter Pitts and Warren McCulloch, who proposed the early mathematical models of neural networks to come up with algorithms that mirror human thought processes.
By the early 1960s, an experimental "learning machine" with punched tape memory, called Cybertron, had been developed by Raytheon Company to analyse sonar signals, electrocardiograms, and speech patterns using rudimentary reinforcement learning. It was repetitively "trained" by a human operator/teacher to recognise patterns and equipped with a "goof" button to cause it to reevaluate incorrect decisions. A representative book on research into machine learning during the 1960s was Nils Nilsson's book on Learning Machines, dealing mostly with machine learning for pattern classification. Interest related to pattern recognition continued into the 1970s, as described by Duda and Hart in 1973. In 1981, a report was given on using teaching strategies so that an artificial neural network learns to recognise 40 characters (26 letters, 10 digits, and 4 special symbols) from a computer terminal.
Tom M. Mitchell provided a widely quoted, more formal definition of the algorithms studied in the machine learning field: "A computer program is said to learn from experience E with respect to some class of tasks T and performance measure P if its performance at tasks in T, as measured by P,  improves with experience E." This definition of the tasks in which machine learning is concerned offers a fundamentally operational definition rather than defining the field in cognitive terms. This follows Alan Turing's proposal in his paper "Computing Machinery and Intelligence", in which the question, "Can machines think?", is replaced with the question, "Can machines do what we (as thinking entities) can do?".
Modern-day Machine Learning algorithms are broken into 3 algorithm types: Supervised Learning Algorithms, Unsupervised Learning Algorithms, and Reinforcement Learning Algorithms.

Current Supervised Learning Algorithms have objectives of classification and regression.
Current Unsupervised Learning Algorithms have objectives of clustering, dimensionality reduction, and association rule.
Current Reinforcement Learning Algorithms focus on decisions that must be made with respect to some previous, unknown time and are broken down to either be studies of model-based methods or model-free methods.
In 2014 Ian Goodfellow and others introduced generative adversarial networks (GANs) with realistic data synthesis. By 2016 AlphaGo obtained victory against top human players using reinforcement learning techniques.


== Relationships to other fields ==


=== Artificial intelligence ===

As a scientific endeavour, machine learning grew out of the quest for artificial intelligence (AI). In the early days of AI as an academic discipline, some researchers were interested in having machines learn from data. They attempted to approach the problem with various symbolic methods, as well as what were then termed "neural networks"; these were mostly perceptrons and other models that were later found to be reinventions of the generalised linear models of statistics. Probabilistic reasoning was also employed, especially in automated medical diagnosis.
However, an increasing emphasis on the logical, knowledge-based approach caused a rift between AI and machine learning. Probabilistic systems were plagued by theoretical and practical problems of data acquisition and representation. By 1980, expert systems had come to dominate AI, and statistics was out of favour. Work on symbolic/knowledge-based learning did continue within AI, leading to inductive logic programming(ILP), but the more statistical line of research was now outside the field of AI proper, in pattern recognition and information retrieval. Neural networks research had been abandoned by AI and computer science around the same time. This line, too, was continued outside the AI/CS field, as "connectionism", by researchers from other disciplines, including John Hopfield, David Rumelhart, and Geoffrey Hinton. Their main success came in the mid-1980s with the reinvention of backpropagation.
Machine learning (ML), reorganised and recognised as its own field, started to flourish in the 1990s. The field changed its goal from achieving artificial intelligence to tackling solvable problems of a practical nature. It shifted focus away from the symbolic approaches it had inherited from AI, and toward methods and models borrowed from statistics, fuzzy logic, and probability theory.


=== Data compression ===


=== Data mining ===
Machine learning and data mining often employ the same methods and overlap significantly, but while machine learning focuses on prediction, based on known properties learned from the training data, data mining focuses on the discovery of (previously) unknown properties in the data (this is the analysis step of knowledge discovery in databases). Data mining uses many machine learning methods, but with different goals; on the other hand, machine learning also employs data mining methods as "unsupervised learning" or as a preprocessing step to improve learner accuracy. Much of the confusion between these two research communities (which do often have separate conferences and separate journals, ECML PKDD being a major exception) comes from the basic assumptions they work with: in machine learning, performance is usually evaluated with respect to the ability to reproduce known knowledge, while in knowledge discovery and data mining (KDD) the key task is the discovery of previously unknown knowledge. Evaluated with respect to known knowledge, an uninformed (unsupervised) method will easily be outperformed by other supervised methods, while in a typical KDD task, supervised methods cannot be used due to the unavailability of training data.
Machine learning also has intimate ties to optimisation: Many learning problems are formulated as minimisation of some loss function on a training set of examples. Loss functions express the discrepancy between the predictions of the model being trained and the actual problem instances (for example, in classification, one wants to assign a label to instances, and models are trained to correctly predict the preassigned labels of a set of examples).


=== Generalization ===
Characterizing the generalisation of various learning algorithms is an active topic of current research, especially for deep learning algorithms.


=== Statistics ===
Machine learning and statistics are closely related fields in terms of methods, but distinct in their principal goal: statistics draws population inferences from a sample, while machine learning finds generalisable predictive patterns.
Conventional statistical analyses require the a priori selection of a model most suitable for the study data set. In addition, only significant or theoretically relevant variables based on previous experience are included for analysis. In contrast, machine learning is not built on a pre-structured model; rather, the data shape the model by detecting underlying patterns. The more variables (input) used to train the model, the more accurate the ultimate model will be.
Leo Breiman distinguished two statistical modelling paradigms: data model and algorithmic model, wherein "algorithmic model" means more or less the machine learning algorithms like Random Forest.
Some statisticians have adopted methods from machine learning, leading to a combined field that they call statistical learning.


=== Statistical physics ===
Analytical and computational techniques derived from deep-rooted physics of disordered systems can be extended to large-scale problems, including machine learning, e.g., to analyse the weight space of deep neural networks. Statistical physics is thus finding applications in the area of medical diagnostics.


== Theory ==

A core objective of a learner is to generalise from its experience. Generalisation in this context is the ability of a learning machine to perform accurately on new, unseen examples/tasks after having experienced a learning data set. The training examples come from some generally unknown probability distribution (considered representative of the space of occurrences) and the learner has to build a general model about this space that enables it to produce sufficiently accurate predictions in new cases.
The computational analysis of machine learning algorithms and their performance is a branch of theoretical computer science known as computational learning theory via the probably approximately correct learning  model. Because training sets are finite and the future is uncertain, learning theory usually does not yield guarantees of the performance of algorithms. Instead, probabilistic bounds on the performance are quite common. The bias–variance decomposition is one way to quantify generalisation error.
For the best performance in the context of generalisation, the complexity of the hypothesis should match the complexity of the function underlying the data. If the hypothesis is less complex than the function, then the model has underfitted the data. If the complexity of the model is increased in response, then the training error decreases. But if the hypothesis is too complex, then the model is subject to overfitting and generalisation will be poorer.
In addition to performance bounds, learning theorists study the time complexity and feasibility of learning. In computational learning theory, a computation is considered feasible if it can be done in polynomial time. There are two kinds of time complexity results: Positive results show that a certain class of functions can be learned in polynomial time. Negative results show that certain classes cannot be learned in polynomial time.


== Approaches ==


